\documentclass[12pt,a4paper]{article}

% ─── Encodage et langue ───────────────────────────────────────────────────────
\usepackage[utf8]{inputenc}
\usepackage[T1]{fontenc}
\usepackage[french]{babel}

% ─── Mise en page ─────────────────────────────────────────────────────────────
\usepackage[top=2.5cm, bottom=2.5cm, left=2.8cm, right=2.8cm]{geometry}
\usepackage{setspace}
\onehalfspacing

% ─── Typographie ──────────────────────────────────────────────────────────────
\usepackage{lmodern}
\usepackage{microtype}

% ─── Couleurs et graphiques ───────────────────────────────────────────────────
\usepackage[dvipsnames,table]{xcolor}
\definecolor{dbtgreen}{HTML}{00855A}
\definecolor{lightgray}{HTML}{F5F5F5}
\definecolor{codebg}{HTML}{F0F0F0}
\definecolor{headercolor}{HTML}{2C3E50}
\definecolor{accentblue}{HTML}{2980B9}
\definecolor{warncolor}{HTML}{E67E22}

% ─── Code source ──────────────────────────────────────────────────────────────
\usepackage{listings}
\lstdefinelanguage{SQL}{
  keywords={SELECT, FROM, WHERE, JOIN, LEFT, RIGHT, INNER, ON, AS,
            WITH, GROUP, BY, ORDER, HAVING, DISTINCT, COUNT, SUM, AVG,
            MIN, MAX, CASE, WHEN, THEN, ELSE, END, INSERT, UPDATE, SET,
            CREATE, TABLE, INDEX, AND, OR, NOT, IN, EXISTS, IS, NULL,
            COALESCE, ROUND, EXTRACT, CURRENT_TIMESTAMP, CURRENT_DATE,
            TRUE, FALSE, INTO, VALUES, DROP, ALTER, ADD, PARTITION,
            OVER, ROW_NUMBER, MD5, CONCAT_WS, PERCENTILE_CONT, WITHIN,
            ANALYZE, TRUNCATE, BETWEEN, ANY},
  keywordstyle=\color{accentblue}\bfseries,
  morestring=[b]',
  stringstyle=\color{Mahogany},
  morecomment=[l]{--},
  commentstyle=\color{gray}\itshape,
  sensitive=false
}
\lstdefinelanguage{Jinja}{
  keywords={config, ref, source, is_incremental, if, endif, macro,
            endmacro, materialized, unique_key, post_hook, on_schema_change},
  keywordstyle=\color{dbtgreen}\bfseries,
  morestring=[b]',
  stringstyle=\color{Mahogany},
  morecomment=[l]{--},
  commentstyle=\color{gray}\itshape,
  sensitive=false
}
\lstset{
  basicstyle=\ttfamily\small,
  backgroundcolor=\color{codebg},
  frame=single,
  rulecolor=\color{gray!40},
  breaklines=true,
  breakatwhitespace=true,
  showstringspaces=false,
  numbers=left,
  numberstyle=\tiny\color{gray},
  numbersep=6pt,
  tabsize=4,
  captionpos=b,
  xleftmargin=1em,
  xrightmargin=0.5em
}

% ─── Tableaux ─────────────────────────────────────────────────────────────────
\usepackage{booktabs}
\usepackage{tabularx}
\usepackage{multirow}
\usepackage{longtable}

% ─── Liens et méta-données ────────────────────────────────────────────────────
\usepackage{hyperref}
\hypersetup{
  colorlinks=true,
  linkcolor=headercolor,
  urlcolor=accentblue,
  citecolor=dbtgreen,
  pdftitle={Rapport SAE -- Pipeline DVF avec historisation SCD2},
  pdfauthor={Belarbi -- Delrive}
}

% ─── En-têtes et pieds de page ────────────────────────────────────────────────
\usepackage{fancyhdr}
\pagestyle{fancy}
\fancyhf{}
\fancyhead[L]{\small\textcolor{gray}{Pipeline DVF -- dbt + SCD2}}
\fancyhead[R]{\small\textcolor{gray}{Belarbi \& Delrive}}
\fancyfoot[C]{\thepage}
\renewcommand{\headrulewidth}{0.4pt}

% ─── Divers ───────────────────────────────────────────────────────────────────
\usepackage{graphicx}
\usepackage{enumitem}
\usepackage{tcolorbox}
\tcbuselibrary{skins,breakable}
\usepackage{titlesec}
\usepackage{caption}
\usepackage{float}

\titleformat{\section}{\Large\bfseries\color{headercolor}}{}{0em}
  {\colorbox{headercolor!10}{\parbox{\dimexpr\linewidth-2\fboxsep}{\thesection\quad #1}}}
\titleformat{\subsection}{\large\bfseries\color{headercolor}}{\thesubsection}{1em}{}
\titleformat{\subsubsection}{\normalsize\bfseries\color{dbtgreen}}{\thesubsubsection}{1em}{}

\newtcolorbox{infobox}[1][]{
  colback=accentblue!5, colframe=accentblue!50,
  title=#1, fonttitle=\bfseries\small, breakable
}
\newtcolorbox{successbox}[1][]{
  colback=dbtgreen!5, colframe=dbtgreen!50,
  title=#1, fonttitle=\bfseries\small, breakable
}
\newtcolorbox{warnbox}[1][]{
  colback=warncolor!5, colframe=warncolor!50,
  title=#1, fonttitle=\bfseries\small, breakable
}

\newcommand{\dbt}{\textbf{\textcolor{dbtgreen}{dbt}}}
\newcommand{\sql}[1]{\texttt{\small #1}}
\newcommand{\schema}[1]{\textbf{\texttt{#1}}}

% ══════════════════════════════════════════════════════════════════════════════
\begin{document}

% ─── Page de titre ────────────────────────────────────────────────────────────
\begin{titlepage}
  \centering
  \vspace*{1cm}
  {\Huge\bfseries\color{headercolor} Rapport de Projet SAE\par}
  \vspace{0.5cm}
  {\large\color{gray}\rule{0.6\linewidth}{1pt}\par}
  \vspace{1cm}
  {\LARGE\bfseries Pipeline de données immobilières\par}
  \vspace{0.3cm}
  {\Large\textcolor{dbtgreen}{Demandes de Valeurs Foncières (DVF)}\par}
  \vspace{0.3cm}
  {\large Entrepôt de données avec historisation complète SCD2\par}
  {\large via \dbt{} et PostgreSQL\par}
  \vspace{1.5cm}
  {\large\color{gray}\rule{0.6\linewidth}{0.5pt}\par}
  \vspace{1cm}
  \begin{tabular}{rl}
    \textbf{Auteurs :}  & Reda Belarbi \& Delrive \\[0.4em]
    \textbf{Projet :}   & \texttt{Data\_Retro\_Belarbi\_Delrive} \\[0.4em]
    \textbf{Stack :}    & dbt 1.11.4 · PostgreSQL · Python/Streamlit \\[0.4em]
    \textbf{Date :}     & Mai 2026 \\
  \end{tabular}
  \vspace{1.5cm}

  \begin{tcolorbox}[colback=dbtgreen!8, colframe=dbtgreen!40, width=0.85\linewidth]
    \centering\small
    Ce rapport décrit l'architecture et l'implémentation complète du pipeline
    ETL/ELT réalisé dans le cadre du projet SAE. Il couvre les quatre couches
    de transformation des données brutes DVF jusqu'à l'entrepôt de données
    analytiques avec historisation SCD2 intégrale, ainsi que le tableau de bord
    interactif de visualisation.
  \end{tcolorbox}
  \vfill
\end{titlepage}

\tableofcontents
\newpage

% ══════════════════════════════════════════════════════════════════════════════
\section{Introduction et contexte}
% ══════════════════════════════════════════════════════════════════════════════

\subsection{Présentation du sujet}

Ce projet porte sur la construction d'un \textbf{pipeline de données complet}
pour l'analyse des mutations immobilières en France, à partir des données
publiques \textbf{DVF} (\textit{Demandes de Valeurs Foncières}). Ces données,
publiées par la Direction Générale des Finances Publiques (DGFiP), recensent
l'ensemble des transactions immobilières (ventes, échanges, adjudications,
expropriations, etc.) réalisées sur le territoire français depuis 2014.

L'objectif est de transformer ces données brutes en un \textbf{entrepôt de
données analytiques} permettant d'effectuer des analyses sur les prix de
l'immobilier, les surfaces, les types de biens et leur localisation
géographique, \textit{tout en conservant l'historique complet des
modifications} grâce à une implémentation rigoureuse du SCD de type 2.

\subsection{Technologies utilisées}

\begin{table}[H]
  \centering
  \begin{tabularx}{\linewidth}{lX}
    \toprule
    \textbf{Technologie} & \textbf{Rôle} \\
    \midrule
    \textbf{PostgreSQL} (OVH VPS) & Base de données relationnelle hébergeant tous les schémas \\
    \dbt{} v1.11.4              & Outil de transformation SQL (Data Build Tool) \\
    \texttt{dbt-postgres} v1.10.0 & Adaptateur PostgreSQL pour dbt \\
    \texttt{dbt\_utils} $\geq$1.0 & Package de macros utilitaires \\
    \textbf{Python 3.13}        & Scripts de visualisation et connexion BDD \\
    \textbf{Streamlit 1.57}     & Framework de tableau de bord interactif \\
    \textbf{Plotly 6.x}         & Bibliothèque de graphiques interactifs \\
    \textbf{Git}                & Versionnement du code \\
    \bottomrule
  \end{tabularx}
  \caption{Technologies employées dans le projet}
\end{table}

\subsection{Architecture globale en couches}

Le pipeline suit une architecture en \textbf{quatre couches successives}, chacune
matérialisée dans un schéma PostgreSQL dédié :

\begin{figure}[H]
  \centering
  \begin{tcolorbox}[colback=gray!5, colframe=gray!30, width=0.9\linewidth]
    \centering\ttfamily\small
    \textbf{[Tables sources DVF brutes]}\\
    $\downarrow$ Lecture par vues\\
    \textcolor{accentblue}{\textbf{Staging}} (schéma \texttt{DSA} -- vues)
    \quad\textit{typage, renommage}\\
    $\downarrow$\\
    \textcolor{accentblue}{\textbf{Intermédiaire}} (schéma \texttt{DSA} -- tables)
    \quad\textit{jointures, agrégations, hash MD5}\\
    $\downarrow$\\
    \textcolor{dbtgreen}{\textbf{ODS}} (schéma \texttt{Intermediate} -- incrémental + SCD2)
    \quad\textit{persistance historisée}\\
    $\downarrow$\\
    \textcolor{Mahogany}{\textbf{DWH Marts}} (schéma \texttt{Datawarehouse} -- étoile + SCD2)
    \quad\textit{analytique}\\
    $\downarrow$\\
    \textbf{Tableau de bord Streamlit/Plotly}
    \quad\textit{visualisation interactive}
  \end{tcolorbox}
  \caption{Architecture en couches du pipeline de données}
\end{figure}

\begin{table}[H]
  \centering
  \rowcolors{2}{gray!10}{white}
  \begin{tabularx}{\linewidth}{llXl}
    \toprule
    \textbf{Couche}       & \textbf{Schéma PG}    & \textbf{Rôle}                          & \textbf{Mat.} \\
    \midrule
    Staging               & \texttt{DSA}          & Typage, renommage, nettoyage           & Vue \\
    Intermédiaire         & \texttt{DSA}          & Jointures, agrégations, hash MD5       & Table \\
    ODS                   & \texttt{Intermediate} & Stockage incrémental + SCD2            & Incrémental \\
    DWH Marts             & \texttt{Datawarehouse}& Dimensions + faits + SCD2 complet     & Incr./Table \\
    \bottomrule
  \end{tabularx}
  \caption{Résumé des couches, schémas PostgreSQL et matérialisations}
\end{table}

% ══════════════════════════════════════════════════════════════════════════════
\section{Configuration du projet dbt}
% ══════════════════════════════════════════════════════════════════════════════

\subsection{Fichier \texttt{dbt\_project.yml}}

\begin{lstlisting}[language=Jinja, caption={dbt\_project.yml -- Configuration principale}]
name: Data_Retro_Belarbi_Delrive
version: '1.0.0'
profile: Data_Retro_Belarbi_Delrive

model-paths: ["models"]
macro-paths:  ["macros"]

models:
  Data_Retro_Belarbi_Delrive:
    staging:
      +materialized: view
      +schema: DSA
    intermediate:
      +materialized: table
      +schema: DSA
    ods:
      +materialized: incremental
      +schema: Intermediate
    marts:
      dimensions:
        +materialized: table
        +schema: Datawarehouse
      facts:
        +materialized: incremental
        +schema: Datawarehouse
\end{lstlisting}

\begin{infobox}[Choix de matérialisation par couche]
\begin{itemize}[noitemsep]
  \item \textbf{Staging $\to$ Vue} : pas de duplication des données sources.
  \item \textbf{Intermédiaire $\to$ Table} : les jointures et agrégations coûteuses sont persistées. Important : ces tables doivent être incluses dans chaque \texttt{dbt run} upstream.
  \item \textbf{ODS $\to$ Incrémental} : seules les nouvelles versions sont insérées, le SCD2 gère l'historique.
  \item \textbf{Marts/faits $\to$ Incrémental} : la table de faits grandit sans retraiter l'historique.
\end{itemize}
\end{infobox}

\subsection{Macro \texttt{generate\_schema\_name} -- Isolation des schémas}

\begin{lstlisting}[language=Jinja, caption={macros/generate\_schema\_name.sql}]
{% macro generate_schema_name(custom_schema_name, node) -%}
    {%- if custom_schema_name is none -%}
        {{ target.schema }}
    {%- else -%}
        {{ custom_schema_name | trim }}
    {%- endif -%}
{%- endmacro %}
\end{lstlisting}

Sans cette surcharge, dbt préfixerait le schéma cible devant le nom de schéma
personnalisé (\texttt{DSA\_DSA} au lieu de \texttt{DSA}). La macro garantit que
les modèles atterrissent exactement dans \texttt{DSA}, \texttt{Intermediate}
et \texttt{Datawarehouse}.

\subsection{Déclaration des sources (\texttt{sources.yml})}

\begin{lstlisting}[language=Jinja, caption={models/staging/sources.yml -- extrait}]
sources:
  - name: DSA
    schema: DSA
    tables:
      - name: mutations
      - name: parcelles
      - name: sous_parcelles
      - name: locals
      - name: lots
      - name: localisations
      - name: communes
      - name: departements
      - name: parcelles_mutations
      - name: nature_mutation
      ...

  - name: Intermediate
    schema: Intermediate
    tables:
      - name: ods_transaction
      - name: ods_bien_immobilier
      ...
\end{lstlisting}

\begin{infobox}[Important : \texttt{source()} vs \texttt{ref()}]
La macro \texttt{source()} ne crée \textbf{pas de dépendance DAG} dans dbt.
Tous les modèles ODS référencés depuis les marts utilisent \texttt{ref()} et
non \texttt{source()} pour garantir l'ordre d'exécution correct. Sans cette
précaution, les marts pourraient s'exécuter avant les tables ODS.
\end{infobox}

% ══════════════════════════════════════════════════════════════════════════════
\section{Couche Staging -- Typage et nettoyage}
% ══════════════════════════════════════════════════════════════════════════════

\subsection{Rôle et principes}

La couche Staging a des responsabilités limitées et précises :
\begin{itemize}[noitemsep]
  \item Caster les types (\sql{::int}, \sql{::date}, \sql{::decimal}, etc.)
  \item Renommer les colonnes si nécessaire
  \item Ajouter un horodatage de chargement (\sql{\_loaded\_at})
  \item \textbf{Aucune logique métier}
\end{itemize}

Tous les modèles staging sont matérialisés en \textbf{vues} pour éviter
la duplication des données sources.

\subsection{Modèle clé : \texttt{stg\_mutations}}

\begin{lstlisting}[language=SQL, caption={models/staging/stg\_mutations.sql}]
select
    id_mutation,
    numero_mutation::int          as numero_mutation,
    date_mutation::date           as date_mutation,
    id_nature_mutation::int       as id_nature_mutation,
    valeur_fonciere::bigint       as valeur_fonciere,
    date_chargement,
    current_timestamp             as _loaded_at
from {{ source('DSA', 'mutations') }}
\end{lstlisting}

La colonne \sql{date\_chargement} est essentielle : elle permet de gérer le
pattern \textbf{append-only} décrit en section~\ref{sec:appendonly}.

\subsection{Récapitulatif des modèles staging}

\begin{table}[H]
  \centering
  \rowcolors{2}{gray!10}{white}
  \begin{tabularx}{\linewidth}{lX}
    \toprule
    \textbf{Modèle} & \textbf{Transformations principales} \\
    \midrule
    \texttt{stg\_mutations}           & Cast types, ajout \sql{date\_chargement} \\
    \texttt{stg\_locals}              & Cast \sql{surface\_batiment::int}, \sql{nombre\_piece\_principal::int} \\
    \texttt{stg\_localisations}       & Cast \sql{latitude/longitude::decimal(10,7)} \\
    \texttt{stg\_lots}                & Cast \sql{lot\_surface\_carrez::decimal(10,2)} \\
    \texttt{stg\_parcelles}           & Renommage \sql{code\_nature\_culture\_speciale} \\
    \texttt{stg\_communes}            & Sélection colonnes, \sql{\_loaded\_at} \\
    \texttt{stg\_departements}        & Sélection colonnes, \sql{\_loaded\_at} \\
    \texttt{stg\_sous\_parcelles}     & Cast \sql{surface\_terrain::decimal(12,2)} \\
    \texttt{stg\_parcelles\_mutations}& Table de liaison mutation $\leftrightarrow$ parcelle \\
    \texttt{stg\_nature\_mutation}    & Cast \sql{id\_nat\_mut::int} \\
    \texttt{stg\_types\_local}        & Cast \sql{id\_type\_loc::int} \\
    \bottomrule
  \end{tabularx}
  \caption{Récapitulatif des modèles staging}
\end{table}

% ══════════════════════════════════════════════════════════════════════════════
\section{Couche Intermédiaire -- Logique métier}
% ══════════════════════════════════════════════════════════════════════════════

\subsection{Rôle}

La couche Intermédiaire applique la logique métier : jointures multi-entités,
agrégations, calcul du hash MD5 pour la détection de changements. Elle produit
des \textbf{tables persistées} servant de base à la couche ODS.

\subsection{\texttt{int\_transactions} -- Gestion du pattern append-only}
\label{sec:appendonly}

La table source \sql{DSA.mutations} est en mode \textbf{append-only} : au lieu
de mettre à jour les lignes existantes, chaque version d'une mutation est
ajoutée comme une nouvelle ligne avec un horodatage \sql{date\_chargement}. Ce
choix architectural permet de conserver toutes les versions successives d'une
mutation dans la source.

\begin{lstlisting}[language=SQL, caption={models/intermediate/int\_transactions.sql}]
with mutations_brutes as (
    select
        m.id_mutation,
        m.numero_mutation,
        m.date_mutation,
        m.id_nature_mutation,
        m.valeur_fonciere,
        pm.id_parcelle,
        row_number() over (
            partition by m.id_mutation
            order by m.date_chargement desc   -- derniere version en premier
        ) as rn
    from {{ ref('stg_mutations') }}           m
    join {{ ref('stg_parcelles_mutations') }} pm
        on pm.id_mutation = m.id_mutation
),

mutations_latest as (
    select * from mutations_brutes where rn = 1   -- une seule ligne par mutation
),

mutations_aggregees as (
    select
        id_mutation, numero_mutation, date_mutation, id_nature_mutation,
        valeur_fonciere,
        count(distinct id_parcelle)   as nb_parcelle,
        min(id_parcelle)              as id_parcelle_principale
    from mutations_latest
    group by id_mutation, numero_mutation, date_mutation,
             id_nature_mutation, valeur_fonciere
)

select
    ma.id_mutation, ma.numero_mutation, ma.date_mutation,
    ma.id_nature_mutation, ma.valeur_fonciere,
    ma.nb_parcelle, ma.id_parcelle_principale,
    sp.id_adresse,
    md5(concat_ws('|',
        ma.id_mutation, ma.date_mutation::text,
        ma.valeur_fonciere::text, ma.id_parcelle_principale
    )) as hash_contenu
from mutations_aggregees             ma
left join {{ ref('stg_sous_parcelles') }} sp
    on sp.id_parcelle = ma.id_parcelle_principale
\end{lstlisting}

\begin{infobox}[Déduplication par \texttt{ROW\_NUMBER}]
La fonction \sql{ROW\_NUMBER() OVER (PARTITION BY id\_mutation ORDER BY
date\_chargement DESC)} sélectionne uniquement la \textbf{version la plus
récente} de chaque mutation depuis la source. Sans cette déduplication,
une mutation ayant deux versions dans \sql{DSA.mutations} produirait des
doublons dans toutes les couches aval.
\end{infobox}

\subsection{\texttt{int\_parcelles} -- Parcelles enrichies}

\begin{lstlisting}[language=SQL, caption={models/intermediate/int\_parcelles.sql (extrait)}]
select
    p.id_parcelle, p.ancien_id_parcelle,
    p.code_nature_culture, p.code_nature_culture_spe,
    sp.surface_terrain,   -- vient de sous_parcelles
    sp.id_adresse,
    md5(concat_ws('|', p.id_parcelle, p.code_nature_culture,
        p.code_nature_culture_spe, sp.surface_terrain::text)) as hash_contenu
from {{ ref('stg_parcelles') }}       p
left join {{ ref('stg_sous_parcelles') }} sp
    on sp.id_parcelle = p.id_parcelle
\end{lstlisting}

\subsection{\texttt{int\_locals} -- Locaux avec comptage des lots}

\begin{lstlisting}[language=SQL, caption={models/intermediate/int\_locals.sql (extrait)}]
with locals_avec_lots as (
    select
        l.id_local, l.id_parcelle, l.code_type_local,
        l.surface_batiment, l.nombre_piece_principal,
        count(lo.id_lot) as nb_lots
    from {{ ref('stg_locals') }}  l
    left join {{ ref('stg_lots') }} lo on lo.id_local = l.id_local
    group by l.id_local, l.id_parcelle, l.code_type_local,
             l.surface_batiment, l.nombre_piece_principal
)
select *,
    md5(concat_ws('|', id_local, code_type_local::text,
        surface_batiment::text, nb_lots::text)) as hash_contenu
from locals_avec_lots
\end{lstlisting}

% ══════════════════════════════════════════════════════════════════════════════
\section{Couche ODS -- Historisation SCD2}
% ══════════════════════════════════════════════════════════════════════════════

\subsection{Principe du SCD Type 2}

Le \textbf{Slowly Changing Dimension de Type 2} (SCD2) conserve
\textit{toutes les versions historiques} d'un enregistrement. Lorsqu'une donnée
change, on n'écrase pas l'ancienne valeur : on \textbf{ferme l'ancienne version}
et on \textbf{insère une nouvelle}.

\begin{table}[H]
  \centering
  \rowcolors{2}{gray!10}{white}
  \begin{tabularx}{\linewidth}{lX}
    \toprule
    \textbf{Colonne SCD2} & \textbf{Signification} \\
    \midrule
    \sql{ods\_sk}         & Clé de substitution = \sql{md5(id || '|' || hash\_contenu)} \\
    \sql{hash\_contenu}   & MD5 des colonnes métier, sert à détecter les changements \\
    \sql{date\_debut}     & Horodatage d'entrée en vigueur de la version \\
    \sql{date\_fin}       & \sql{'9999-12-31'} si courant, sinon date de fermeture \\
    \sql{est\_courant}    & \sql{true} pour la version active, \sql{false} pour les archives \\
    \sql{statut\_*}       & \sql{'VALIDE'} ou \sql{'EXPIRE'} selon le cycle de vie \\
    \bottomrule
  \end{tabularx}
  \caption{Colonnes de gestion SCD2 dans les tables ODS}
\end{table}

\subsection{Macro \texttt{scd2\_close\_records}}

\begin{lstlisting}[language=Jinja, caption={macros/scd2\_close\_records.sql}]
{% macro scd2_close_records(target_table, id_col, hash_col,
                             statut_col='statut') %}
UPDATE {{ target_table }} AS cible
SET
    date_fin         = current_timestamp,
    est_courant      = false,
    {{ statut_col }} = 'EXPIRE'
WHERE cible.est_courant = true
AND EXISTS (
    SELECT 1 FROM {{ target_table }} nouvelle
    WHERE nouvelle.{{ id_col }}   = cible.{{ id_col }}
    AND   nouvelle.{{ hash_col }} != cible.{{ hash_col }}
    AND   nouvelle.est_courant    = true
    AND   nouvelle.date_debut     > cible.date_debut
);
{% endmacro %}
\end{lstlisting}

\begin{successbox}[Fonctionnement en deux temps]
\begin{enumerate}[noitemsep]
  \item \textbf{INSERT} : la nouvelle version est insérée avec
  \sql{est\_courant = true} et \sql{date\_fin = '9999-12-31'}.
  \item \textbf{UPDATE} (post\_hook) : \sql{scd2\_close\_records} ferme
  l'ancienne version en positionnant \sql{est\_courant = false},
  \sql{date\_fin = now()} et \sql{statut = 'EXPIRE'}.
\end{enumerate}
La macro n'agit que si deux versions coexistent avec des hash différents
(\sql{hash\_contenu != ...}), évitant les fermetures parasites.
\end{successbox}

\subsection{\texttt{ods\_transaction} -- Transactions avec SCD2 complet}

\begin{lstlisting}[language=Jinja, caption={models/ods/ods\_transaction.sql}]
{{
  config(
    materialized = 'incremental',
    unique_key   = 'ods_sk',
    post_hook    = [
      "{{ scd2_close_records(this, 'id_mutation', 'hash_contenu',
                             'statut_transaction') }}",
      "CREATE INDEX IF NOT EXISTS idx_ods_transaction_est_courant
           ON {{ this }} (est_courant)",
      "CREATE INDEX IF NOT EXISTS idx_ods_transaction_hash
           ON {{ this }} (hash_contenu) WHERE est_courant = true",
      "ANALYZE {{ this }}"
    ]
  )
}}

with source as (select * from {{ ref('int_transactions') }}),

nouveaux as (
    select
        md5(concat_ws('|', s.id_mutation, s.hash_contenu)) as ods_sk,
        gen_random_uuid()                   as id_transaction,
        s.id_mutation,
        s.id_parcelle_principale            as id_parcelle,
        s.id_nature_mutation,
        s.numero_mutation, s.date_mutation, s.valeur_fonciere,
        s.nb_parcelle, s.id_adresse, s.hash_contenu,
        current_timestamp                   as date_debut,
        '9999-12-31 00:00:00'::timestamp    as date_fin,
        true                                as est_courant,
        'VALIDE'                            as statut_transaction,
        current_timestamp                   as date_creation_ods,
        1                                   as version
    from source s
    {% if is_incremental() %}
    where not exists (
        select 1 from {{ this }} t
        where t.hash_contenu = s.hash_contenu and t.est_courant = true
    )
    {% endif %}
)
select * from nouveaux
\end{lstlisting}

\begin{infobox}[Clé de substitution \sql{ods\_sk}]
\sql{ods\_sk = md5(id\_mutation || '|' || hash\_contenu)} garantit l'unicité
par \textbf{version} : deux versions d'une même mutation ont le même
\sql{id\_mutation} mais des \sql{ods\_sk} différents. C'est cette clé qui est
utilisée comme \sql{unique\_key} pour le chargement incrémental dbt.
\end{infobox}

\subsection{\texttt{ods\_parcelle} et \texttt{ods\_bien\_immobilier}}

Ces modèles suivent le même pattern SCD2 : calcul d'un \sql{ods\_sk},
filtre incrémental sur \sql{hash\_contenu}, et fermeture des anciens
enregistrements via \sql{scd2\_close\_records} en post\_hook.

\texttt{ods\_bien\_immobilier} est le modèle le plus complexe de la couche
ODS : il consolide les informations de parcelle et de local en une entité
unique de bien immobilier, permettant ainsi d'agréger les surfaces et le
nombre de lots au niveau de la mutation dans le DWH.

\subsection{Récapitulatif des modèles ODS}

\begin{table}[H]
  \centering
  \rowcolors{2}{gray!10}{white}
  \begin{tabularx}{\linewidth}{llX}
    \toprule
    \textbf{Modèle} & \textbf{SCD2} & \textbf{Particularité} \\
    \midrule
    \texttt{ods\_transaction}      & Oui & Hash sur id\_mutation + valeur\_fonciere + date \\
    \texttt{ods\_parcelle}         & Oui & Hash sur nature culture + surface terrain \\
    \texttt{ods\_local}            & Oui & Hash sur surface batiment + type + nb lots \\
    \texttt{ods\_bien\_immobilier} & Oui & Consolidation parcelle + local \\
    \texttt{ods\_localisation}     & Oui & Hash sur adresse + code commune \\
    \texttt{ods\_commune}          & Oui & Hash sur code département + code postal \\
    \texttt{ods\_departement}      & Oui & Hash sur libellé département \\
    \texttt{ods\_ref\_nature\_mutation} & Oui & Référentiel des types de mutation \\
    \bottomrule
  \end{tabularx}
  \caption{Modèles ODS et implémentation SCD2}
\end{table}

% ══════════════════════════════════════════════════════════════════════════════
\section{Couche DWH Marts -- Entrepôt analytique avec SCD2 intégral}
% ══════════════════════════════════════════════════════════════════════════════

\subsection{Architecture en étoile avec historisation complète}

Le DWH implémente un \textbf{schéma en étoile} autour de la table de faits
\sql{fact\_mutation}. L'originalité de ce projet est que \textbf{toutes les
versions historiques} des faits sont conservées, pas seulement la version
courante -- ce qui va au-delà d'un DWH classique.

\begin{figure}[H]
  \centering
  \begin{tcolorbox}[colback=gray!5, colframe=gray!30, width=0.85\linewidth]
    \centering\ttfamily\small
    \textbf{dim\_temps} (id\_date)\\
    \quad $\searrow$\\
    \quad\quad \textbf{fact\_mutation} $\longleftarrow$ \textbf{localisation} (id\_location)\\
    \quad $\nearrow$\\
    \textbf{type\_mutation} (id\_type\_mutation)\\[0.3em]
    \textit{Les trois tables portent des colonnes SCD2 :\\
    sk (surrogate key), date\_debut, date\_fin, est\_courant}
  \end{tcolorbox}
  \caption{Schéma en étoile du DWH avec SCD2 intégral}
\end{figure}

\subsection{\texttt{dim\_temps} -- Dimension temporelle}

\begin{lstlisting}[language=SQL, caption={models/marts/dimensions/dim\_temps.sql}]
{{ config(materialized='table') }}

with spine as (
    select generate_series(
        '2000-01-01'::date, current_date, '1 day'::interval
    )::date as date_
)
select
    date_,
    extract(year    from date_)::int        as annee,
    extract(month   from date_)::int        as mois,
    to_char(date_, 'TMMonth')               as lib_mois,
    extract(quarter from date_)::int        as trimestre,
    'T' || extract(quarter from date_)::int as lib_trimestre,
    extract(dow from date_)::int            as jour_semaine
from spine
\end{lstlisting}

La fonction \sql{generate\_series} PostgreSQL génère automatiquement toutes
les dates depuis le 1\textsuperscript{er} janvier 2000 jusqu'à aujourd'hui.

\subsection{\texttt{type\_mutation} -- Dimension avec SCD2}

\begin{lstlisting}[language=Jinja, caption={models/marts/dimensions/type\_mutation.sql}]
{{
  config(
    materialized = 'incremental',
    unique_key   = 'sk_type_mutation'
  )
}}

select
    ref.ods_sk                  as sk_type_mutation,
    ref.id_nat_mut::varchar     as id_type_mutation,
    ref.lib_nat_mut             as libelle_type_mutation,
    ref.date_debut,
    ref.date_fin,
    ref.est_courant
from {{ ref('ods_ref_nature_mutation') }} ref

{% if is_incremental() %}
where not exists (
    select 1 from {{ this }} t
    where t.sk_type_mutation = ref.ods_sk
)
{% endif %}
\end{lstlisting}

\subsection{\texttt{localisation} -- Dimension géographique avec SCD2}

\begin{lstlisting}[language=Jinja, caption={models/marts/dimensions/localisation.sql}]
{{
  config(
    materialized = 'incremental',
    unique_key   = 'sk_localisation'
  )
}}

select
    loc.ods_sk                  as sk_localisation,
    loc.id_adresse::varchar     as id_localisation,
    d.code_dep,
    d.lib_dep                   as lib_departement,
    left(d.code_dep, 2)         as reg,
    d.lib_dep                   as lib_reg,
    c.code_com,
    c.lib_com,
    c.code_postal,
    loc.date_debut,
    loc.date_fin,
    loc.est_courant
from {{ ref('ods_localisation') }}  loc
join {{ ref('ods_commune') }}        c
    on  c.code_com    = loc.code_com
    and c.est_courant = true
join {{ ref('ods_departement') }}    d
    on  d.code_dep    = c.code_dep
    and d.est_courant = true

{% if is_incremental() %}
where not exists (
    select 1 from {{ this }} t
    where t.sk_localisation = loc.ods_sk
)
{% endif %}
\end{lstlisting}

\begin{infobox}[Propagation du SCD2 dans les dimensions]
Les dimensions \sql{localisation} et \sql{type\_mutation} propagent les
colonnes SCD2 (\sql{date\_debut}, \sql{date\_fin}, \sql{est\_courant}) depuis
l'ODS. Cela permet d'interroger l'état d'une dimension \textbf{à une date
précise} : \sql{WHERE est\_courant = true} donne l'état actuel ;
une condition sur \sql{date\_debut}/\sql{date\_fin} donne l'état historique.
\end{infobox}

\subsection{\texttt{fact\_mutation} -- Table de faits avec historisation complète}

C'est la pièce centrale du DWH. Sa particularité est de conserver
\textbf{toutes les versions historiques} d'une mutation (pas seulement la
version courante), et de synchroniser son état avec l'ODS via un
\sql{post\_hook}.

\begin{lstlisting}[language=Jinja, caption={models/marts/facts/fact\_mutation.sql}]
{{
  config(
    materialized     = 'incremental',
    unique_key       = 'sk_mutation',
    on_schema_change = 'sync_all_columns',
    post_hook        = ["""
      UPDATE {{ this }} fm
      SET   est_courant = ot.est_courant,
            date_fin    = ot.date_fin
      FROM  {{ ref('ods_transaction') }} ot
      WHERE fm.id_mutation  = ot.id_mutation
        AND fm.sk_mutation  = ot.ods_sk
        AND fm.est_courant != ot.est_courant
    """]
  )
}}

with transactions as (
    select t.*
    from {{ ref('ods_transaction') }} t
    where t.statut_transaction in ('VALIDE', 'EXPIRE')  -- toutes les versions

    {% if is_incremental() %}
    and not exists (  -- n'insere pas ce qui est deja dans fact_mutation
        select 1 from {{ this }} fm
        where fm.sk_mutation = t.ods_sk
    )
    {% endif %}
),

biens_par_mutation as (
    select
        t.id_mutation,
        avg(b.surface_batiment)   as avg_surface_bat,
        avg(b.surface_terrain)    as avg_surface_terrain,
        avg(b.nb_lots)            as avg_nb_local
    from transactions t
    left join {{ ref('ods_bien_immobilier') }} b
        on  b.id_parcelle = t.id_parcelle
        and b.est_courant = true
    group by t.id_mutation
)

select
    t.ods_sk                                    as sk_mutation,
    t.id_mutation,
    t.id_nature_mutation::varchar               as id_type_mutation,
    t.date_mutation                             as id_date,
    t.id_adresse::varchar                       as id_location,
    t.valeur_fonciere,
    round(bm.avg_surface_bat::numeric, 2)       as avg_surface_bat,
    round(bm.avg_surface_terrain::numeric, 2)   as avg_surface_terrain,
    t.nb_parcelle,
    round(bm.avg_nb_local::numeric, 2)          as avg_nb_local,
    t.date_debut,
    t.date_fin,
    t.est_courant,
    t.date_creation_ods
from transactions     t
join biens_par_mutation bm
    on  bm.id_mutation      = t.id_mutation
join {{ ref('dim_temps') }}     dt
    on  dt.date_            = t.date_mutation
join {{ ref('type_mutation') }} tm
    on  tm.id_type_mutation = t.id_nature_mutation::varchar
    and tm.est_courant      = true
\end{lstlisting}

\begin{table}[H]
  \centering
  \rowcolors{2}{gray!10}{white}
  \begin{tabularx}{\linewidth}{llX}
    \toprule
    \textbf{Colonne} & \textbf{Type} & \textbf{Description} \\
    \midrule
    \sql{sk\_mutation}         & varchar & Clé de substitution = \sql{ods\_sk} de l'ODS \\
    \sql{id\_mutation}         & varchar & Identifiant métier de la mutation \\
    \sql{id\_type\_mutation}   & varchar & FK $\to$ \sql{type\_mutation} \\
    \sql{id\_date}             & date    & FK $\to$ \sql{dim\_temps} \\
    \sql{id\_location}         & varchar & FK $\to$ \sql{localisation} \\
    \sql{valeur\_fonciere}     & bigint  & Prix de vente en euros \\
    \sql{avg\_surface\_bat}    & numeric & Surface bâtie moyenne (m\textsuperscript{2}) \\
    \sql{avg\_surface\_terrain}& numeric & Surface du terrain moyen (m\textsuperscript{2}) \\
    \sql{nb\_parcelle}         & int     & Nombre de parcelles dans la mutation \\
    \sql{avg\_nb\_local}       & numeric & Nombre moyen de lots \\
    \sql{date\_debut}          & timestamp & Début de validité de la version \\
    \sql{date\_fin}            & timestamp & Fin de validité (\sql{'9999-12-31'} si courant) \\
    \sql{est\_courant}         & bool    & \sql{true} = version active \\
    \sql{date\_creation\_ods}  & timestamp & Horodatage de chargement ODS \\
    \bottomrule
  \end{tabularx}
  \caption{Colonnes de \texttt{fact\_mutation}}
\end{table}

\begin{successbox}[Synchronisation SCD2 via post\_hook]
Le \sql{post\_hook} de \sql{fact\_mutation} est un \sql{UPDATE} qui propage
les changements d'état depuis l'ODS. Quand \sql{scd2\_close\_records} ferme
une version dans \sql{ods\_transaction} (\sql{est\_courant} passe à
\sql{false}), le post\_hook détecte le décalage (\sql{fm.est\_courant !=
ot.est\_courant}) et met à jour \sql{fact\_mutation} en conséquence.
Ainsi la table de faits reste toujours cohérente avec l'ODS sans retraitement
complet.
\end{successbox}

% ══════════════════════════════════════════════════════════════════════════════
\section{Tests de qualité des données}
% ══════════════════════════════════════════════════════════════════════════════

\subsection{Tests génériques dbt (\texttt{schema.yml})}

\begin{lstlisting}[language=Jinja, caption={models/marts/schema.yml -- extrait}]
models:
  - name: localisation
    columns:
      - name: sk_localisation
        tests: [unique, not_null]
      - name: est_courant
        tests:
          - accepted_values:
              values: [true, false]
      - name: date_debut
        tests: [not_null]

  - name: fact_mutation
    columns:
      - name: sk_mutation
        tests: [unique, not_null]
      - name: valeur_fonciere
        tests: [not_null]
      - name: est_courant
        tests:
          - accepted_values:
              values: [true, false]
\end{lstlisting}

\subsection{Tests SCD2 personnalisés}

Trois tests SQL personnalisés vérifient la cohérence de l'historisation.
Ils retournent des lignes en cas d'échec (convention dbt).

\subsubsection{Test 1 -- Unicité de la version courante (\texttt{scd2\_courant\_unique.sql})}

\begin{lstlisting}[language=SQL, caption={tests/scd2\_courant\_unique.sql}]
-- Echec si une cle naturelle a plus d'un enregistrement courant
select 'localisation' as modele,
       id_localisation as cle_naturelle,
       count(*) as nb_courants
from {{ ref('localisation') }}
where est_courant = true
group by id_localisation
having count(*) > 1

union all

select 'type_mutation', id_type_mutation, count(*)
from {{ ref('type_mutation') }}
where est_courant = true
group by id_type_mutation
having count(*) > 1
\end{lstlisting}

\subsubsection{Test 2 -- Cohérence des dates (\texttt{scd2\_dates\_coherentes.sql})}

\begin{lstlisting}[language=SQL, caption={tests/scd2\_dates\_coherentes.sql}]
-- Echec si : courant avec date_fin != 9999-12-31,
--            ou clos avec date_fin = 9999-12-31,
--            ou date_debut >= date_fin
select 'localisation', sk_localisation::text,
       est_courant, date_debut, date_fin
from {{ ref('localisation') }}
where (est_courant = true  and date_fin::date != '9999-12-31')
   or (est_courant = false and date_fin::date  = '9999-12-31')
   or date_debut >= date_fin
\end{lstlisting}

\subsubsection{Test 3 -- Absence de chevauchement (\texttt{scd2\_pas\_chevauchement.sql})}

\begin{lstlisting}[language=SQL, caption={tests/scd2\_pas\_chevauchement.sql}]
-- Echec si deux versions d'une meme cle naturelle se chevauchent
select 'localisation', a.id_localisation,
       a.sk_localisation::text, b.sk_localisation::text
from {{ ref('localisation') }} a
join {{ ref('localisation') }} b
    on  a.id_localisation  = b.id_localisation
    and a.sk_localisation != b.sk_localisation
    and a.date_debut        < b.date_fin
    and a.date_fin          > b.date_debut

union all

select 'type_mutation', a.id_type_mutation,
       a.sk_type_mutation::text, b.sk_type_mutation::text
from {{ ref('type_mutation') }} a
join {{ ref('type_mutation') }} b
    on  a.id_type_mutation  = b.id_type_mutation
    and a.sk_type_mutation != b.sk_type_mutation
    and a.date_debut         < b.date_fin
    and a.date_fin           > b.date_debut
\end{lstlisting}

\begin{table}[H]
  \centering
  \rowcolors{2}{gray!10}{white}
  \begin{tabularx}{\linewidth}{lX}
    \toprule
    \textbf{Test} & \textbf{Invariant vérifié} \\
    \midrule
    \texttt{scd2\_courant\_unique}      & Au plus 1 version \sql{est\_courant=true} par clé naturelle \\
    \texttt{scd2\_dates\_coherentes}    & Cohérence \sql{date\_debut}/\sql{date\_fin}/\sql{est\_courant} \\
    \texttt{scd2\_pas\_chevauchement}   & Pas de chevauchement de périodes de validité \\
    Generic \sql{unique}                & Unicité des surrogate keys (\sql{sk\_*}) \\
    Generic \sql{not\_null}             & Colonnes obligatoires non nulles \\
    Generic \sql{accepted\_values}      & \sql{est\_courant} $\in \{true, false\}$ \\
    \bottomrule
  \end{tabularx}
  \caption{Récapitulatif des tests de qualité}
\end{table}

% ══════════════════════════════════════════════════════════════════════════════
\section{Flux de données complet}
% ══════════════════════════════════════════════════════════════════════════════

\begin{longtable}{lllll}
  \toprule
  \textbf{Modèle} & \textbf{Couche} & \textbf{Schéma} & \textbf{Mat.} & \textbf{SCD2} \\
  \midrule
  \endfirsthead
  \toprule
  \textbf{Modèle} & \textbf{Couche} & \textbf{Schéma} & \textbf{Mat.} & \textbf{SCD2} \\
  \midrule
  \endhead
  \midrule\multicolumn{5}{r}{\small (suite...)}\\
  \endfoot
  \bottomrule
  \endlastfoot
  stg\_mutations            & Staging & DSA           & Vue    & Non \\
  stg\_parcelles            & Staging & DSA           & Vue    & Non \\
  stg\_sous\_parcelles      & Staging & DSA           & Vue    & Non \\
  stg\_locals               & Staging & DSA           & Vue    & Non \\
  stg\_lots                 & Staging & DSA           & Vue    & Non \\
  stg\_localisations        & Staging & DSA           & Vue    & Non \\
  stg\_communes             & Staging & DSA           & Vue    & Non \\
  stg\_departements         & Staging & DSA           & Vue    & Non \\
  stg\_parcelles\_mutations & Staging & DSA           & Vue    & Non \\
  stg\_nature\_mutation     & Staging & DSA           & Vue    & Non \\
  \midrule
  int\_transactions         & Interm. & DSA           & Table  & Non \\
  int\_parcelles            & Interm. & DSA           & Table  & Non \\
  int\_locals               & Interm. & DSA           & Table  & Non \\
  int\_localisations        & Interm. & DSA           & Table  & Non \\
  \midrule
  ods\_transaction          & ODS & Intermediate & Incr. & \textbf{Oui} \\
  ods\_parcelle             & ODS & Intermediate & Incr. & \textbf{Oui} \\
  ods\_local                & ODS & Intermediate & Incr. & \textbf{Oui} \\
  ods\_bien\_immobilier     & ODS & Intermediate & Incr. & \textbf{Oui} \\
  ods\_localisation         & ODS & Intermediate & Incr. & \textbf{Oui} \\
  ods\_commune              & ODS & Intermediate & Incr. & \textbf{Oui} \\
  ods\_departement          & ODS & Intermediate & Incr. & \textbf{Oui} \\
  ods\_ref\_nature\_mutation& ODS & Intermediate & Incr. & \textbf{Oui} \\
  \midrule
  dim\_temps                & DWH & Datawarehouse & Table  & Non \\
  type\_mutation            & DWH & Datawarehouse & Incr.  & \textbf{Oui} \\
  localisation              & DWH & Datawarehouse & Incr.  & \textbf{Oui} \\
  fact\_mutation            & DWH & Datawarehouse & Incr.  & \textbf{Oui} \\
\end{longtable}

% ══════════════════════════════════════════════════════════════════════════════
\section{Tableau de bord interactif}
% ══════════════════════════════════════════════════════════════════════════════

\subsection{Architecture technique}

Un tableau de bord interactif a été développé en Python avec
\textbf{Streamlit} et \textbf{Plotly}, connecté directement au schéma
\sql{Datawarehouse} via \textbf{psycopg2}.

\begin{infobox}[Optimisation des performances]
Toutes les requêtes SQL \textbf{agrègent les données côté serveur PostgreSQL}.
On ne rapatrie jamais les 7 millions de lignes brutes de \sql{fact\_mutation}
en Python -- seulement les résultats agrégés (quelques centaines de lignes).
Des techniques spécifiques sont utilisées :
\begin{itemize}[noitemsep]
  \item \sql{width\_bucket()} pour l'histogramme (60 buckets calculés en SQL)
  \item \sql{PERCENTILE\_CONT()} pour les quartiles (Q1/médiane/Q3 en SQL)
  \item \sql{DATE\_TRUNC()} pour les agrégations temporelles
  \item \sql{@st.cache\_data(ttl=120)} pour mettre en cache les résultats
\end{itemize}
\end{infobox}

\subsection{Fonctionnalités}

\textbf{Filtres globaux (sidebar) :}
\begin{itemize}[noitemsep]
  \item Plage de dates (date\_input avec bornes min/max depuis la BDD)
  \item Type de mutation (multiselect)
  \item Département (multiselect)
  \item Tranche de valeur foncière (slider)
\end{itemize}

\textbf{KPIs dynamiques :} nombre de mutations, volume total (Md€),
valeur moyenne (k€), valeur médiane (k€), nombre de communes.

\textbf{5 onglets d'analyse :}

\begin{table}[H]
  \centering
  \rowcolors{2}{gray!10}{white}
  \begin{tabularx}{\linewidth}{lX}
    \toprule
    \textbf{Onglet} & \textbf{Visualisations} \\
    \midrule
    Temporel      & Barres (nb mutations), courbe (valeur moy.), area (volume), heatmap mois$\times$année \\
    Géographique  & Top N départements/communes, treemap département$\to$commune \\
    Par type      & Donut (répartition), barres (valeur moy.), area empilée temporelle \\
    Distributions & Histogramme (buckets SQL), box plot quartiles, scatter surface/valeur, barres nb parcelles \\
    SCD2          & Timeline Gantt des versions, courbe d'évolution de valeur, tableau détaillé \\
    \bottomrule
  \end{tabularx}
  \caption{Onglets et visualisations du tableau de bord}
\end{table}

\begin{lstlisting}[language=SQL, caption={Exemple : histogramme via width\_bucket PostgreSQL}]
SELECT width_bucket(fm.valeur_fonciere, v_min, v_max, 60) AS bucket,
       COUNT(*)           AS nb,
       MIN(fm.valeur_fonciere) AS val_min,
       MAX(fm.valeur_fonciere) AS val_max
FROM "Datawarehouse"."fact_mutation" fm
JOIN "Datawarehouse"."type_mutation" tm ...
WHERE fm.est_courant = true AND ...
GROUP BY bucket ORDER BY bucket
\end{lstlisting}

% ══════════════════════════════════════════════════════════════════════════════
\section{Difficultés rencontrées et solutions}
% ══════════════════════════════════════════════════════════════════════════════

\begin{table}[H]
  \centering
  \rowcolors{2}{gray!10}{white}
  \begin{tabularx}{\linewidth}{lX}
    \toprule
    \textbf{Problème} & \textbf{Solution apportée} \\
    \midrule
    \sql{UnicodeDecodeError 'charmap' 0x90} & Suppression d'un caractère \texttt{UTF-8} invalide (\texttt{$\leftarrow$}) dans \texttt{dbt\_project.yml} \\
    Marts s'exécutent avant l'ODS & Remplacement de \sql{source()} par \sql{ref()} dans les marts pour créer les dépendances DAG \\
    Disque plein sur le VPS & Suppression des anciens schémas (\texttt{DROP SCHEMA "ODS" CASCADE}) \\
    Contrainte PK sur \texttt{DSA.mutations} & \texttt{ALTER TABLE ... DROP CONSTRAINT mutations\_pkey} pour le mode append-only \\
    Colonne \texttt{date\_chargement} manquante & \texttt{ALTER TABLE ... ADD COLUMN IF NOT EXISTS date\_chargement} \\
    \texttt{fact\_mutation} ne capturait que 1 version & Filtre \sql{statut\_transaction = 'VALIDE'} excluait les versions expirées (\texttt{'EXPIRE'}) -- corrigé en \sql{IN ('VALIDE', 'EXPIRE')} \\
    Tables intermédiaires non rafraîchies & \texttt{dbt run --select int\_transactions+} pour inclure tous les upstream \\
    \texttt{Decimal} PostgreSQL incompatible pandas & Conversion avec \sql{try: pd.to\_numeric(col) except: pass} dans \texttt{run\_query} \\
    \texttt{errors=''ignore''} supprimé en pandas 3.x & Migration vers \sql{try/except} autour de \sql{pd.to\_numeric} \\
    \bottomrule
  \end{tabularx}
  \caption{Difficultés rencontrées et solutions}
\end{table}

% ══════════════════════════════════════════════════════════════════════════════
\section{Conclusion}
% ══════════════════════════════════════════════════════════════════════════════

\subsection{Bilan technique}

Ce projet implémente un pipeline de données complet, robuste et historisé
pour les données DVF françaises. Les principales réalisations sont :

\begin{enumerate}
  \item \textbf{Architecture 4 couches} avec séparation stricte des
  responsabilités : Staging (typage), Intermédiaire (logique métier),
  ODS (historisation), DWH (analytique).

  \item \textbf{SCD Type 2 généralisé} : implémenté sur 8 tables ODS
  et propagé jusqu'au DWH (dimensions et table de faits), permettant
  d'analyser l'état du marché immobilier à \textit{n'importe quelle date}.

  \item \textbf{Pattern append-only} : la source DVF peut recevoir plusieurs
  versions d'une même mutation ; la déduplication par \sql{ROW\_NUMBER()
  OVER (PARTITION BY id\_mutation ORDER BY date\_chargement DESC)} isole
  la version la plus récente.

  \item \textbf{Synchronisation post\_hook} : \sql{fact\_mutation} se
  synchronise automatiquement avec l'ODS sans retraitement complet grâce
  à un \sql{UPDATE} ciblé en post\_hook.

  \item \textbf{Tests SCD2 personnalisés} : trois tests SQL garantissent
  l'intégrité de l'historique (unicité du courant, cohérence des dates,
  absence de chevauchement).

  \item \textbf{Tableau de bord interactif} : 5 axes d'analyse, filtres
  dynamiques, toutes les requêtes agrégées côté PostgreSQL pour des
  performances optimales sur 7+ millions de lignes.
\end{enumerate}

\subsection{Volumes de données}

\begin{table}[H]
  \centering
  \rowcolors{2}{gray!10}{white}
  \begin{tabularx}{\linewidth}{lrl}
    \toprule
    \textbf{Table} & \textbf{Lignes} & \textbf{Remarque} \\
    \midrule
    \texttt{fact\_mutation}   & $\sim$7,2 M & Toutes versions SCD2 \\
    \texttt{localisation}     & $\sim$5,3 M & Toutes versions SCD2 \\
    \texttt{type\_mutation}   & 8            & Référentiel compact \\
    \texttt{dim\_temps}       & $\sim$9 000  & 2000 $\to$ aujourd'hui \\
    \bottomrule
  \end{tabularx}
  \caption{Volumes de données dans le schéma Datawarehouse}
\end{table}

\subsection{Perspectives}

\begin{itemize}[noitemsep]
  \item Intégration d'un orchestrateur (Airflow, Dagster) pour automatiser
  les runs dbt quotidiens en réponse à de nouveaux fichiers DVF.
  \item Enrichissement avec des données INSEE (revenus par commune,
  démographie) pour des analyses socio-économiques.
  \item Ajout d'une couche de géolocalisation (PostGIS) pour des
  visualisations cartographiques.
  \item Extension du tableau de bord avec une carte choroplèthe par
  département.
\end{itemize}

\end{document}
